\section{总结与展望}

本文围绕“利用 VGGT 的几何能力进行三维编辑与生成”开展文献调研、模型复现和系统设计。通过 DUSt3R、VGGT、FLARE、VGGT-Edit、RL3DEdit、TRELLIS 和 Vinedresser3D 等工作可以看到，当前三维内容创建正在从多阶段逐场景优化转向前馈基础模型、原生三维潜空间和多模型协同。

本文已完成 VGGT 的多视图几何输出与可视化验证，并对 VGGT 与 TRELLIS 1 的官方代码进行了接口级核对。代码证据表明，VGGT 能够输出相机、深度、pointmap 与置信度，并导出点云 GLB 或 COLMAP 数据；TRELLIS 具备图像/文本条件、两阶段采样和 mesh/3DGS/Radiance Field 多格式解码能力。归档的 \texttt{room.ply} 确实包含 3,551,987 个完整 Gaussian 参数，但仓库没有保存从当前 VGGT 输出训练到该 PLY 的脚本与日志；原先归因于 TRELLIS 的 \texttt{classroom.glb} 经结构解析实为彩色点云与相机线框。

在明确这些实现边界后，本文提出以 VGGT 负责真实场景几何、TRELLIS 负责新增三维内容、相似变换与支撑面约束负责空间对齐的融合方案，并设计了语言驱动定位、局部高斯编辑、多视图一致性、背景保持和空间合理性评价。该方案属于可执行的工程设计，尚不是当前仓库已跑通的端到端系统。

后续应首先补齐可复现证据链：保存实际运行环境、输入帧列表、命令行、权重版本、TRELLIS condition/seed/采样参数和逐阶段耗时；然后实现生成对象添加到教室场景的最小闭环，比较直接插入、支撑面对齐、完整几何对齐和 VGGT 几何筛选四种设置。在此基础上，可从以下方向扩展：

\begin{itemize}[itemsep=0.15em, topsep=0.35em]
  \item 使用 SAM、REALM 或 VLM 将自然语言指令转换为跨视图一致的三维 mask；
  \item 借鉴 VGGT-Edit，在 mask 内预测 Gaussian 参数的局部残差；
  \item 借鉴 RL3DEdit，将 VGGT 几何反馈从后验评分扩展为训练奖励；
  \item 使用 TRELLIS 2 的 O-Voxel 与 PBR 材质增强拓扑和重光照能力；
  \item 面向动态场景，引入 Easi3R 或 VGGT-$\Omega$ 的动静分离与动态重建。
\end{itemize}

总体而言，VGGT 可以作为连接观察、生成和编辑的统一几何接口；与 TRELLIS 结合后，有望以较低训练成本实现可解释、可交互且多视图一致的真实场景编辑。
